\chapter{Procedencia de los datos}
\label{chap:anexoB}

\noindent Este anexo tabula el origen de todas las series usadas en la investigación empírica
(Capítulo~\ref{chap:paper2}) y en la prueba de las predicciones distintivas del modelo
(Capítulo~\ref{chap:paper1}, Sección~\ref{sec:datos_reales}). El panel se construye sobre el
principio de \textbf{homogeneidad de fuente}: todos los insumos de mercado provienen de la
terminal Bloomberg, extraídos con un mismo protocolo para las veinte economías del universo
objetivo; solo los componentes reales del índice de condiciones financieras que alimenta
$GaR$ y los controles macroeconómicos domésticos provienen de estadísticas oficiales, dado
que Bloomberg no es fuente primaria de cuentas nacionales. El código de extracción y
consolidación reside en \texttt{1\_Codigo/Bloomberg\_extraction/} y \texttt{1\_Codigo/Panel/bbg/}.

\vspace{1em}

\begin{footnotesize}
\setlength{\tabcolsep}{3.5pt}
\renewcommand{\arraystretch}{1.15}
\newcommand{\Lc}{\raggedright\arraybackslash\hyphenpenalty=50\relax\hspace{0pt}}
\begin{longtable}{%
  >{\Lc}p{2.15cm}%
  >{\Lc}p{2.7cm}%
  >{\Lc}p{1.8cm}%
  >{\Lc}p{2.85cm}%
  >{\Lc}p{1.75cm}%
  >{\Lc}p{2.5cm}}
\caption{Procedencia, definición y transformación de cada serie del panel.}
\label{tab:datos}\\
\toprule
\textbf{Serie} & \textbf{Definición} & \textbf{Fuente} & \textbf{Identificador} & \textbf{Frec.\ nativa} & \textbf{Transformación} \\
\midrule
\endfirsthead
\multicolumn{6}{l}{\footnotesize\itshape Tabla~\ref{tab:datos} (continuación)}\\
\toprule
\textbf{Serie} & \textbf{Definición} & \textbf{Fuente} & \textbf{Identificador} & \textbf{Frec.\ nativa} & \textbf{Transformación} \\
\midrule
\endhead
\midrule
\multicolumn{6}{r}{\footnotesize\itshape continúa en la página siguiente}\\
\endfoot
\bottomrule
\endlastfoot

\multicolumn{6}{l}{\textbf{\textit{Variable dependiente}}}\\
\addlinespace
EMBI Global Diversified & Diferencial del índice de bonos soberanos de mercados emergentes de J.P.\ Morgan frente a la curva del Tesoro estadounidense (\textbf{variable dependiente principal}, siguiendo a \citealp{Chari2024}) & J.P.\ Morgan (serie \texttt{embi.xlsx}) & spread EMBI Global por país & Diaria & Media trimestral, puntos básicos. Sin dato $\Rightarrow$ la observación no entra a la regresión de referencia \\
CDS soberano 5Y & Prima del \textit{credit default swap} soberano a cinco años, en dólares, senior no garantizado (\textbf{serie de robustez}) & Bloomberg & \texttt{<emisor> CDS USD SR 5Y Corp}; curva Markit \texttt{CCHIN1U5 Curncy} para China & Diaria & Media trimestral, puntos básicos. Correlación 0,86 con el EMBI en la submuestra común \\
\addlinespace
\multicolumn{6}{l}{\textbf{\textit{Fragilidad bancaria --- $JLoss$}} (métrica compuesta; motor \texttt{jloss\_engine.py})}\\
\addlinespace
Activos totales & Activos totales del banco & Bloomberg & \texttt{BS\_TOT\_ASSET} & Trimestral & Por banco cotizado; ajuste de periodicidad \texttt{ACTUAL} \\
Patrimonio & Patrimonio total contable & Bloomberg & \texttt{TOTAL\_EQUITY} & Trimestral & --- \\
Deuda emitida (LP) & \textit{Long-term borrowings} (bonos, letras, deuda subordinada); define el largo plazo del punto de incumplimiento & Bloomberg & \texttt{BS\_LT\_BORROW} & Trimestral & $D^{*}=$ (pasivo total $-$ bonos) $+ \tfrac12$ bonos \\
Ingreso neto, ingresos, caja & Insumos del $z$-score y de controles de coherencia contable & Bloomberg & \texttt{NET\_\allowbreak INCOME}, \texttt{SALES\_\allowbreak REV\_\allowbreak TURN}, \texttt{BS\_\allowbreak CASH\_\allowbreak NEAR\_\allowbreak CASH\_\allowbreak ITEM} & Trimestral & --- \\
Capital. bursátil & Valor de mercado del patrimonio, cotización local (nunca ADR), verificada contra \texttt{EQY\_\allowbreak FUND\_\allowbreak CRNCY} & Bloomberg & \texttt{CUR\_\allowbreak MKT\_\allowbreak CAP}, o \texttt{PX\_\allowbreak LAST} $\times$ \texttt{BS\_\allowbreak SH\_\allowbreak OUT} & Diaria & Último valor del trimestre \\
Retornos accionarios & Retornos logarítmicos diarios del precio local & Bloomberg & \texttt{PX\_LAST} & Diaria & $\sigma_E=\sqrt{\sum_t r_t^2}$ por banco-trimestre (varianza realizada) \\
Índice bursátil país & Factor sistémico único de la agregación de pérdidas & Bloomberg & índice de referencia por país (\texttt{COUNTRY\_INDEX}) & Diaria & Retornos $\to$ correlación $\rho_i$ banco-factor \\
\addlinespace
\multicolumn{6}{p{15cm}}{\footnotesize\textit{Parámetros del motor (idénticos a \citealp{Chari2024}):} LGD $=0{,}45$; $\rho=0{,}4$; 7 nodos de cuadratura Gauss--Hermite; factor sistémico $N(0,\,0{,}16\%)$; percentil 99; cota de pérdidas $[0{,}010,\,0{,}048]$; horizonte 1 trimestre.}\\
\addlinespace
\multicolumn{6}{l}{\textbf{\textit{Riesgo de cola --- $GaR$}} (regresión cuantílica de panel; motor \texttt{gar\_engine.py}, \texttt{fci\_engine.py})}\\
\addlinespace
PIB real & Crecimiento trimestral del producto (variable dependiente de la regresión cuantílica) & Estadística nacional (banco central) & \texttt{GDP\_*} & Trimestral & Tasa de variación \\
IPC & Componente del FCI (bloque de deuda / expectativas) & Estadística nacional & \texttt{CPI\_*} & Mensual & Estand.\ expansiva \\
Índice accionario & Componente del FCI (bloque accionario) & Bolsa nacional & \texttt{STX\_*} & Mensual & Estand.\ expansiva \\
Tipo de cambio real ef. & Componente del FCI (bloque cambiario) & BIS / fuentes nacionales & \texttt{rEER\_*} & Mensual & Estand.\ expansiva \\
Rendimiento soberano 10Y & Componente del FCI (bloque de deuda soberana) & Mercado / estadística nacional & \texttt{Ryr\_*} & Mensual & Estand.\ expansiva \\
VIX (ortogonalizado) & Condiciones financieras globales en la regresión cuantílica & Bloomberg & \texttt{VIX Index} & Diaria & Ortogonalizado respecto de las condiciones locales; trimestral \\
\addlinespace
\multicolumn{6}{l}{\textbf{\textit{Factores financieros globales}}}\\
\addlinespace
VIX & Índice de volatilidad implícita S\&P 500 & Bloomberg & \texttt{VIX Index} & Diaria & Media trimestral \\
Tasa del Tesoro 10Y & Rendimiento del bono del Tesoro estadounidense a 10 años & Bloomberg & \texttt{USGG10YR Index} & Diaria & Media trimestral; en logaritmo \\
\textit{Spread} HY EE.UU. & \textit{Option-adjusted spread} del \textit{high yield} corporativo estadounidense (equivale a FRED \texttt{BAMLH0A0HYM2}) & Bloomberg & \texttt{LF98OAS Index} & Diaria & Media trimestral; en logaritmo \\
\textit{Spread} on/off-run & Diferencial de liquidez del Tesoro (instrumento \textit{shift-share} principal) & Serie previa del proyecto & \texttt{global\_controls\_} \texttt{quarterly.csv} & Trimestral & En logaritmo \\
Dólar efectivo amplio & Tipo de cambio efectivo nominal de EE.\ UU., canasta de 64 economías (instrumento \textit{shift-share} secundario) & BIS, \textit{Effective Exchange Rate Indices} & \texttt{WS\_\allowbreak EER}, \texttt{M.N.B.US} & Mensual & Media trimestral; en logaritmo \\
\addlinespace
\multicolumn{6}{l}{\textbf{\textit{Controles macroeconómicos domésticos}} (reconstruidos para todos los países del panel)}\\
\addlinespace
Deuda gob.\ general/PIB & Deuda bruta del gobierno general, \% del PIB & FMI, \textit{World Economic Outlook} & \texttt{GGXWDG\_NGDP} & Anual & Interpolación lineal a trimestral \\
Balance fiscal/PIB & Endeudamiento neto del gobierno general, \% del PIB & FMI, \textit{World Economic Outlook} & \texttt{GGXCNL\_NGDP} & Anual & Interpolación lineal a trimestral \\
Reservas/PIB & Reservas internacionales totales sobre PIB corriente & Banco Mundial & \texttt{FI.\allowbreak RES.\allowbreak TOTL.\allowbreak CD} $\div$ \texttt{NY.\allowbreak GDP.\allowbreak MKTP.\allowbreak CD} & Anual & Interpolación lineal a trimestral \\
Cuenta corriente/PIB & Saldo de cuenta corriente, \% del PIB & Banco Mundial & \texttt{BN.\allowbreak CAB.\allowbreak XOKA.\allowbreak GD.\allowbreak ZS} & Anual & Interpolación lineal a trimestral \\
Inflación interanual & Variación del IPC a 12 meses & Serie de IPC del FCI; respaldo Banco Mundial & \texttt{FP.\allowbreak CPI.\allowbreak TOTL.\allowbreak ZG} & Mensual & Variación a 4 trimestres \\
Tipo de cambio real ef. & Índice de tipo de cambio real efectivo (2010 = 100) & BIS; series locales & \texttt{rEER\_*} & Mensual & Último valor del trimestre \\
\addlinespace
\multicolumn{6}{l}{\textbf{\textit{Concentración bancaria e instituciones}}}\\
\addlinespace
$HHI$ (concentración 3 bancos) & Activos de los tres mayores bancos comerciales sobre activos del sistema, \% & Banco Mundial, \textit{Global Financial Development Database} & \texttt{GFDD.\allowbreak OI.01} & Anual & Nivel estructural $=$ mediana por país; serie anual como robustez \\
$HHI_q$ (concentración trimestral) & Activos de los tres mayores bancos \emph{cotizados} sobre activos cotizados totales, \% (mismos bancos que $\JLoss$) & Bloomberg (balances, \texttt{tot\_\allowbreak asset}) & motor \texttt{p6\_\allowbreak concentracion\_\allowbreak trimestral.py} & Trimestral & CR3 y HHI directos; NaN si $<$3 bancos cotizados; corr.\ $\approx$0,5 con GFDD \\
Rating soberano S\&P & Calificación de largo plazo en moneda extranjera & Bloomberg & \texttt{RTG\_\allowbreak SP\_\allowbreak LT\_\allowbreak FC\_\allowbreak ISSUER\_\allowbreak CREDIT} & --- & \textit{Snapshot} (no serie); escala 21--1 \\
Calidad de gobierno (WGI) & Índice de \textit{Worldwide Governance Indicators} (triple interacción institucional) & Banco Mundial & \texttt{instituciones.\allowbreak csv} (provisional) & Anual & Estandarizado; marcado provisional \\

\end{longtable}
\end{footnotesize}

\vspace{1em}

\noindent \textbf{Economías del panel y su cobertura.} El panel incluye toda economía
con $JLoss$ disponible. La muestra de estimación del mecanismo (observaciones con EMBI $+$
$JLoss$ $+$ $GaR$ simultáneamente) reúne \textbf{trece países}: nueve con EMBI de historia
continua (Brasil, Chile, China, Colombia, Malasia, México, Perú, Polonia, Turquía), India
con 54 trimestres (EMBI desde 2012Q4), y Filipinas, Indonesia y Sudáfrica con once cada uno
(EMBI desde 2023Q3, aunque su CDS es de historia larga). Argentina, Egipto y Rusia tienen
$JLoss$ pero no $GaR$; Pakistán tiene $JLoss$ y $GaR$ pero no EMBI. Corea del Sur, Bulgaria
y Hungría se excluyen porque su $JLoss$ no es una medida válida a nivel de país —número de
bancos cotizados insuficiente— (Capítulo~\ref{chap:paper2}, Sección~5.1). La
Figura~\ref{fig:cobertura} traza esta cobertura.

\vspace{1ex}

\noindent La Tabla~\ref{tab:diagpais} resume, por país de la muestra de estimación, el número
mediano de bancos cotizados que alimentan $JLoss$, los trimestres marcados con menos bancos
que el mínimo requerido, y la mediana y el máximo de la serie de $JLoss$. Deja ver dos rasgos
que el Capítulo~\ref{chap:paper2} discute: (i) nueve de los trece países tienen un $JLoss$
mediano entre 2,2 y 2,6 —la variación transversal de la fragilidad la aportan Turquía, India,
China y Brasil—; y (ii) Colombia y México descansan sobre pocas instituciones en una fracción
de los trimestres, aunque sin llegar al caso de Bulgaria (un banco) o Hungría (dos, todos los
trimestres bajo el mínimo) que motivaron su exclusión del panel.

\begin{table}[htbp]
\centering
\footnotesize
\caption{Diagnóstico de $JLoss$ por país (muestra de estimación).}
\label{tab:diagpais}
\begin{tabular}{lrrrr}
\toprule
País & Bancos (mediana) & Trim.\ bajo mínimo & $JLoss$ mediana & $JLoss$ máx. \\
\midrule
Brasil       & 8  & 0  & 4,0  & 15,5 \\
Chile        & 4  & 0  & 2,2  & 11,6 \\
China        & 9  & 2  & 3,8  & 29,0 \\
Colombia     & 3  & 23 & 2,4  & 16,7 \\
Filipinas    & 5  & 0  & 2,9  & 4,3  \\
India        & 11 & 6  & 7,5  & 23,2 \\
Indonesia    & 4  & 0  & 2,3  & 3,7  \\
Malasia      & 6  & 0  & 2,2  & 5,9  \\
México       & 4  & 13 & 2,4  & 13,9 \\
Perú         & 3  & 6  & 2,2  & 8,7  \\
Polonia      & 8  & 0  & 2,6  & 15,0 \\
Sudáfrica    & 6  & 0  & 2,4  & 3,9  \\
Turquía      & 8  & 0  & 6,0  & 28,0 \\
\bottomrule
\end{tabular}

\medskip
\footnotesize\raggedright Fuente: \texttt{1\_Codigo/Panel/bbg/diag\_por\_pais\_bbg.csv}
(motor \texttt{jloss\_engine.py}). ``Trim.\ bajo mínimo'' cuenta los trimestres con menos
bancos cotizados que el umbral del motor. Corea del Sur, Bulgaria y Hungría no aparecen:
quedan fuera del panel por $JLoss$ no válido a nivel de país.
\end{table}
